% sections/related-work.tex

\section{Related Work}
\label{sec:related}

\paragraph{LLM-as-a-Judge.}
The use of LLMs as automated evaluators has grown rapidly.
\citet{zheng2023judging} introduced Chatbot Arena and MT-Bench, establishing pairwise LLM judging as a scalable evaluation paradigm.
\citet{li2024llmsasjudges} provided a comprehensive survey of LLM-based evaluation methods, identifying biases including position bias, length bias, and self-enhancement bias.
\citet{wang2023faireval} proposed mitigating position bias via multi-turn debate and swap-and-compare.
\citet{wei2024systematic} systematically evaluated LLM-as-a-judge prompt sensitivity, finding that rubric template diversity affects alignment scores---a finding that motivates our use of rubric perturbation as an error signal rather than merely a nuisance to be controlled.

\paragraph{Selective Prediction and Abstention.}
Selective classification~\citep{geifman2017selective,el2010foundations} enables a predictor to abstain on uncertain inputs.
\citet{wen2025know} surveyed abstention in LLMs, covering verbalized confidence, self-consistency, and retrieval-based approaches.
\citet{krishnan2024enhancing} proposed uncertainty-calibrated fine-tuning for LLM trust, using last-layer hidden states.
Our work differs by focusing on \emph{evaluation-protocol stability} as the uncertainty source, which is specific to the judging task and does not require model internals.

\paragraph{Risk-Controlled Evaluation.}
\citet{jung2025trust} (Trust-or-Escalate) proposed cascaded selective evaluation for LLM judges with provable human-agreement guarantees via Simulated Annotators and fixed-sequence testing.
Our framework builds on their risk-control machinery but replaces the single-signal confidence (simulated annotator agreement) with a multi-signal protocol-stability calibrator.
We additionally introduce rubric perturbation as a novel signal and provide a direct comparison to their approach.

\paragraph{Judge Benchmarks.}
\citet{tan2024judgebench} introduced JudgeBench, a benchmark of challenging response pairs with objective correctness labels across knowledge, reasoning, math, and coding.
\citet{lambert2024rewardbench} proposed RewardBench for evaluating reward models.
We use JudgeBench as a hard-objective benchmark and TL;DR~\citep{stiennon2020learning} and MT-Bench Human~\citep{zheng2023judging} as subjective-preference benchmarks, covering both correctness and preference evaluation.

\paragraph{Rubric-Based Evaluation.}
\citet{hong2026from} studied aligning LLM judges with human scoring rubrics.
\citet{anghel2025pearl} proposed a rubric-driven multi-metric evaluation framework.
These works treat rubrics as fixed evaluation criteria to be aligned with; in contrast, we treat \emph{rubric sensitivity}---the degree to which a judge's verdict changes under semantically equivalent rubric paraphrases---as an uncertainty signal.
To our knowledge, no prior work has used rubric perturbation instability as a predictor of judge error.
